\subsection{مبانی نظری و فرمول‌بندی ریاضی}
فرض کنید $\mathbf{X} \in \mathbb{R}^{B \times L \times D}$ ماتریس فعال‌سازی ورودی لایه ترنسفورمر باشد. در این معماری، ورودی به طور هم‌زمان وارد شاخه توجه و شاخه بازگشتی \lr{Mamba-2} می‌گردد:

\paragraph{۱. شاخه توجه متراکم مجهز به لورا:}
پروجکشن‌های ورودی با ماتریس‌های کم‌مرتبه تطبیقی فاکتورگیری شده و سپس توجه چندسره استاندارد محاسبه می‌شود:
\begin{equation}
\mathbf{Y}_{\text{attn}} = \operatorname{Softmax}\left(\frac{\mathbf{Q} \mathbf{K}^T}{\sqrt{d}} + \mathbf{M}\right) \mathbf{V} \mathbf{W}_O
\end{equation}

\paragraph{۲. شاخه مدل حالت فضای Mamba-2:}
در شاخه \lr{Mamba2}، تانسور ورودی ابتدا به یک فضای دوبرابری نگاشت شده و به دو بخش تبدیل حالت و گیت مدولاسیون تفکیک می‌شود: $[\mathbf{X}_{\text{ssm}}, \mathbf{Z}] = \operatorname{in\_proj}(\mathbf{X})$. سپس کانولوشن یک‌بعدی زمانی بر روی $\mathbf{X}_{\text{ssm}}$ اعمال شده و پارامترهای پیوسته پیاده‌سازی سیستم خطی به صورت گسسته وابسته به گام زمانی تبدیل می‌شوند:
\begin{align}
\Delta_t &= \operatorname{Softplus}\left(\mathbf{W}_{\Delta} \mathbf{x}_{\text{ssm}, t}\right), \\
\mathbf{A} &= -\exp\left(\operatorname{clamp}(\mathbf{A}_{\text{log}}, -20, 10)\right), \\
\bar{\mathbf{A}}_t &= \exp\left(\operatorname{clamp}(\Delta_t \mathbf{A}, -20, 0)\right), \quad \bar{\mathbf{B}}_t = \mathbf{x}_{\text{ssm}, t} \otimes \Delta_t
\end{align}
رابطه بازگشتی خطی حالت فضا در طول دنباله زمانی با پایداری عددی محاسبه می‌گردد:
\begin{equation}
\mathbf{h}_t = \bar{\mathbf{A}}_t \odot \mathbf{h}_{t-1} + \bar{\mathbf{B}}_t
\end{equation}
\begin{equation}
\mathbf{Y}_{\text{ssm}, t} = \operatorname{out\_proj}\left( \left(\mathbf{h}_t \odot \mathbf{D}\right) \odot \operatorname{SiLU}(\mathbf{z}_t) \right)
\end{equation}

\paragraph{۳. تلفیق دینامیک شاخه‌ها و شبکه خبره‌ها (MoE):}
برآیند این دو شاخه با یک پارامتر یادگیرنده گیتینگ $\alpha$ ترکیب خطی می‌شود:
\begin{equation}
g = \sigma(\alpha), \quad \tilde{\mathbf{X}} = g \cdot \mathbf{Y}_{\text{attn}} + (1 - g) \cdot \mathbf{Y}_{\text{ssm}}
\end{equation}
در نهایت، شبکه خبره‌ها با ساختار برگزیدن برترین‌ها \lr{(Top-$k$ Routing)} بر روی $\tilde{\mathbf{X}}$ اعمال می‌گردد:
\begin{equation}
\mathbf{H}_{\text{moe}} = \sum_{e \in \operatorname{Top-}k(\mathbf{R}(\tilde{\mathbf{X}}))} \operatorname{Softmax}(\mathbf{R}(\tilde{\mathbf{X}}))_e \cdot \operatorname{Expert}_e(\tilde{\mathbf{X}})
\end{equation}
\begin{equation}
\mathbf{X}_{\text{final}} = \tilde{\mathbf{X}} + \gamma_{\text{moe}} \cdot \mathbf{H}_{\text{moe}}
\end{equation}


\paragraph{تحلیل طیفی و نامساوی‌های پایداری خطا}
برای بررسی میزان دقت تقریب ماتریس وزن یا نگاشت توجه در ماژول \lr{granite\_hybrid}، تفاضل ماتریس فاکتورگیری‌شده از نگاشت متراکم اصلی بر اساس نامساوی هافدینگ-فون‌نیومان ارزیابی می‌گردد. فرض کنید $\mathbf{W} \in \mathbb{R}^{d_{\text{in}} \times d_{\text{out}}}$ ماتریس وزن هدف باشد که توسط تقریب کم‌مرتبه یا پراکنده با رتبه موثر $r$ جایگزین شده است. کران بالای خطای بازسازی بر مبنای نرم فروبنیوس به صورت زیر تصریح می‌شود:
\begin{equation}
\label{eq:spectral_bound_granite_hybrid}
\| \mathbf{W} - \mathbf{W}_{r} \|_F = \sqrt{\sum_{i=r+1}^{\min(d_{\text{in}}, d_{\text{out}})} \sigma_i^2(\mathbf{W})} \le \epsilon \| \mathbf{W} \|_F
\end{equation}
که در آن $\sigma_i(\mathbf{W})$ بیانگر $i$-امین مقدار منفرد مرتب‌شده ماتریس به صورت نزولی است. نسبت حفظ انرژی طیفی $\eta_r = \frac{\sum_{i=1}^r \sigma_i^2}{\sum_{i=1}^{d} \sigma_i^2}$ تضمین می‌کند که با انتخاب آستانه $\tau \ge 0.90$، بیش از نود درصد از پراکندگی واریانس داده‌ها حفظ شده و افت سیگنال در انتشار گرادیان رو به عقب به حداقل ممکن کاهش می‌یابد.

\paragraph{تحلیل مجانبی و پیچیدگی اشتراک حافظه سخت‌افزاری}
از منظر پیچیدگی محاسباتی و انتقال داده‌ها در سلسومراتب حافظه (حافظه سراسری \lr{HBM} به حافظه محلی \lr{SRAM})، عملیات ضرب ماتریسی در این مکانیزم نیازمند تخصیص بافر موقت با ابعاد فاکتورگیری‌شده است. برای پردازش دنباله‌ای به طول $N$ با بعد مخفی $D$، بار محاسباتی هسته از فرم درجه دو به رابطه زیر تقلیل می‌یابد:
\begin{equation}
\label{eq:complexity_granite_hybrid}
\mathcal{C}_{\text{granite\_hybrid}} = \mathcal{O}\left( B \cdot N \cdot D \cdot r + B \cdot H \cdot N \cdot d_k \right) \ll \mathcal{O}\left( B \cdot H \cdot N^2 \cdot d_k \right)
\end{equation}
این بهبود مرتبه پیچیدگی، با کاهش چشمگیر حجم تبادل تانسورها در گذرگاه‌های ارتباطی کارت گرافیک، مانع از اشباع پهنای باند حافظه (پدیده \lr{Memory-Bound}) در فرآیند استنتاج طولانی‌مدت می‌گردد.
